\documentclass[11pt,a4paper]{article}

\usepackage[T1]{fontenc}
\usepackage[utf8]{inputenc}
\usepackage{lmodern}
\usepackage[margin=2.3cm]{geometry}
\usepackage{amsmath,amssymb}
\usepackage{booktabs}
\usepackage{enumitem}
\usepackage{microtype}
\usepackage[hidelinks]{hyperref}

\setlength{\parindent}{0pt}
\setlength{\parskip}{0.5em}
\setlist{itemsep=2pt,topsep=4pt}

\newcommand{\dd}{\mathrm{d}}
\newcommand{\J}{\mathcal{J}}

\title{Calculus of Variations:\\How to Differentiate a Curve}
\author{Lecture notes and instructor calculations}
\date{}

\begin{document}

\maketitle

\textbf{Duration:} 45 minutes.

\textbf{Background:} Differentiation, integration by parts, and basic
partial derivatives.

\textbf{Scope:} One independent variable and one unknown function.

\textbf{Central idea:}
Ordinary calculus, but the variable is a whole curve.

\begin{center}
\begin{tabular}{@{}ll@{}}
\toprule
Time & Topic \\
\midrule
0--5 min   & The brachistochrone challenge \\
5--11 min  & From functions to integral functionals \\
11--18 min & What does ``vary the curve'' mean? \\
18--30 min & Derive the Euler--Lagrange equation \\
30--39 min & Solve a complete example \\
39--43 min & Recover Newton's second law \\
43--45 min & Three questions to finish \\
\bottomrule
\end{tabular}
\end{center}

\section*{0--5 minutes: Begin with a race}

Draw two points on the board: \(A\) above and to the left of \(B\).
Draw a straight track and a curved track that drops steeply at first.

\begin{quote}
A bead starts from rest and slides without friction.
Which track gets it to \(B\) first?
\end{quote}

Give the audience a moment to vote.

The straight line has the shortest distance. But a steeper initial
descent builds speed earlier. A longer journey can take less time
if more of it happens at high speed.

In June 1696, Johann Bernoulli challenged mathematicians to find the
curve of fastest descent---the \emph{brachistochrone}.
Solutions published in 1697 included those of Johann and Jacob
Bernoulli, Leibniz, and Newton. The answer is an arc of a cycloid.
A historical reference is given at the end of these notes.

Ask:
\begin{quote}
Do the beads have different speeds when they reach \(B\)?
\end{quote}

They have the \textbf{same final speed}, because they have lost the
same gravitational potential energy. Their travel times differ
because their speeds differ along the way.

Transition:
\begin{quote}
The unknown here is an entire curve.
How do we differentiate with respect to that?
\end{quote}

\section*{5--11 minutes: A functional takes a function and returns a number}

Ordinary calculus deals with maps such as
\[
f:\mathbb{R}\longrightarrow\mathbb{R},
\qquad
z\longmapsto f(z).
\]

A functional can take a whole function as its input:
\[
\J:\{\text{admissible functions}\}\longrightarrow\mathbb{R},
\qquad
y\longmapsto\J[y].
\]

The output could be a length, travel time, energy, or action.

The form we will study is
\[
\boxed{
\J[y]
=
\int_a^b F\bigl(x,y(x),y'(x)\bigr)\,\dd x.
}
\]

Explain the three arguments:
\begin{itemize}
\item \(x\): where we are;
\item \(y(x)\): the value of the curve there;
\item \(y'(x)\): its slope there.
\end{itemize}

The integrand assigns a local contribution.
The integral adds those contributions along the interval.

\subsection*{Constructing the brachistochrone functional}

Take
\[
A=(0,0),
\qquad
B=(b,h),
\qquad
b,h>0,
\]
and measure \(y\) \textbf{vertically downward}.

Assume uniform gravity, no friction, and a bead released from rest.
For a track represented as \(y=y(x)\), take \(y(x)>0\) for \(x>0\).

Conservation of energy gives
\[
\frac12 mv^2=mgy,
\]
so
\[
v(y)=\sqrt{2gy}.
\]

An element of distance along the track is
\[
\dd s
=
\sqrt{\dd x^2+\dd y^2}
=
\sqrt{1+\left(\frac{\dd y}{\dd x}\right)^2}\,\dd x
=
\sqrt{1+y'^2}\,\dd x.
\]

Since \(\dd t=\dd s/v\),
\[
\dd t
=
\frac{\sqrt{1+y'^2}}{\sqrt{2gy}}\,\dd x.
\]

Therefore,
\[
\boxed{
T[y]
=
\int_0^b
\frac{\sqrt{1+y'(x)^2}}{\sqrt{2g\,y(x)}}\,\dd x,
\qquad
y(0)=0,\quad y(b)=h.
}
\]

Point to the numerator and denominator:
\begin{quote}
A steep track adds distance, but getting lower increases speed.
This integral keeps both accounts.
\end{quote}

We now have the quantity to optimize.
What remains is learning how to differentiate it.

\textit{Instructor detail.}
The integral is improper at the release point \(y=0\).
Use it for motivation; derive the general equation first under
ordinary smoothness assumptions.

\section*{11--18 minutes: Turn a wiggle into an ordinary variable}

Fix the endpoint values:
\[
y(a)=\alpha,
\qquad
y(b)=\beta.
\]

Choose a smooth function \(\eta(x)\) satisfying
\[
\eta(a)=\eta(b)=0.
\]

Define a family of nearby curves:
\[
\boxed{
y_\varepsilon(x)=y(x)+\varepsilon\eta(x).
}
\]

Here:
\begin{itemize}
\item \(\eta(x)\) specifies the \textbf{shape} of the wiggle;
\item \(\varepsilon\) specifies its \textbf{size and sign}.
\end{itemize}

Draw the original curve and two nearby curves, all sharing the
same endpoints.

The endpoint conditions are preserved:
\[
y_\varepsilon(a)
=
y(a)+\varepsilon\eta(a)
=
\alpha,
\]
and similarly,
\[
y_\varepsilon(b)=\beta.
\]

Differentiating with respect to \(x\),
\[
y_\varepsilon'(x)
=
y'(x)+\varepsilon\eta'(x).
\]

Emphasize:
\begin{quote}
Changing a curve also changes its slope.
Those two changes are linked.
\end{quote}

For a fixed choice of \(\eta\), define the ordinary scalar function
\[
\Phi_\eta(\varepsilon)
=
\J[y+\varepsilon\eta].
\]

If \(y\) is a local minimum of \(\J\), then \(\varepsilon=0\)
is a local minimum of \(\Phi_\eta\). Thus,
\[
\Phi_\eta'(0)=0.
\]

This derivative is called the \textbf{first variation}:
\[
\boxed{
\delta\J[y;\eta]
:=
\left.
\frac{\dd}{\dd\varepsilon}
\J[y+\varepsilon\eta]
\right|_{\varepsilon=0}.
}
\]

The stationarity requirement is
\[
\boxed{
\delta\J[y;\eta]=0
\quad
\text{for every admissible }\eta.
}
\]

Underline \textbf{every}.

Testing one wiggle tests only one direction in the space of curves.
A minimum must resist all admissible small changes.

\subsection*{Optional analogy with gradients}

For an ordinary function of several variables,
\[
\left.
\frac{\dd}{\dd\varepsilon}
f(\mathbf{z}+\varepsilon\mathbf{v})
\right|_{\varepsilon=0}
=
\nabla f(\mathbf{z})\cdot\mathbf{v}.
\]

Requiring this to vanish for every \(\mathbf{v}\) gives
\[
\nabla f(\mathbf{z})=0.
\]

We are about to find the corresponding statement for functions.

\section*{18--30 minutes: Derive Euler--Lagrange}

Start by substituting the perturbed curve:
\[
\J[y_\varepsilon]
=
\int_a^b
F\bigl(x,y+\varepsilon\eta,y'+\varepsilon\eta'\bigr)\,\dd x.
\]

The independent variable \(x\), and the integration limits, stay fixed.

Assuming sufficient smoothness, differentiate under the integral
and apply the chain rule:
\[
\frac{\dd}{\dd\varepsilon}\J[y_\varepsilon]
=
\int_a^b
\left[
F_y\bigl(x,y_\varepsilon,y_\varepsilon'\bigr)\eta
+
F_{y'}\bigl(x,y_\varepsilon,y_\varepsilon'\bigr)\eta'
\right]\dd x.
\]

Evaluate at \(\varepsilon=0\):
\[
\boxed{
\delta\J[y;\eta]
=
\int_a^b
\left(
F_y\,\eta+F_{y'}\,\eta'
\right)\dd x.
}
\]

From here onward, derivatives of \(F\) are evaluated along the
original curve:
\[
\bigl(x,y(x),y'(x)\bigr).
\]

\subsection*{Clarifying the partial derivatives}

Temporarily view the integrand as \(F(x,u,p)\), with three
independent arguments.

Take partial derivatives with respect to \(u\) and \(p\) first,
then substitute
\[
u=y(x),
\qquad
p=y'(x).
\]

The notation \(F_y\) and \(F_{y'}\) refers to differentiation with
respect to the second and third arguments, respectively.

\subsection*{Integration by parts}

The problem is now the \(\eta'\) term.
We want everything multiplying the same arbitrary function \(\eta\).

Use integration by parts:
\[
\int_a^b F_{y'}\,\eta'\,\dd x
=
\left[F_{y'}\eta\right]_a^b
-
\int_a^b
\frac{\dd}{\dd x}\left(F_{y'}\right)\eta\,\dd x.
\]

Consequently,
\[
\delta\J[y;\eta]
=
\left[F_{y'}\eta\right]_a^b
+
\int_a^b
\left[
F_y-\frac{\dd}{\dd x}\left(F_{y'}\right)
\right]\eta\,\dd x.
\]

Because the endpoints are fixed,
\[
\eta(a)=\eta(b)=0,
\]
and hence
\[
\left[F_{y'}\eta\right]_a^b=0.
\]

We are left with
\[
\boxed{
\delta\J[y;\eta]
=
\int_a^b
\left[
F_y-\frac{\dd}{\dd x}\left(F_{y'}\right)
\right]\eta\,\dd x.
}
\]

Pause:
\begin{quote}
We moved a derivative off the arbitrary wiggle and onto something
determined by the curve.
\end{quote}

\subsection*{Why the coefficient must vanish}

Let
\[
G(x)
=
F_y-\frac{\dd}{\dd x}\left(F_{y'}\right).
\]

Stationarity says
\[
\int_a^b G(x)\eta(x)\,\dd x=0
\]
for every smooth \(\eta\) vanishing at the endpoints.

Why must \(G=0\)?

Suppose \(G(x_0)>0\) somewhere in the interior.
If \(G\) is continuous, it remains positive on a small interval
around \(x_0\).

Choose a nonnegative, nonzero smooth bump \(\eta\) supported
entirely inside that interval. Then
\[
\int_a^b G(x)\eta(x)\,\dd x>0,
\]
contradicting stationarity.

The same argument excludes \(G(x_0)<0\).

This is the \textbf{fundamental lemma of the calculus of variations}.
It gives
\[
\boxed{
\frac{\partial F}{\partial y}
-
\frac{\dd}{\dd x}
\left(
\frac{\partial F}{\partial y'}
\right)
=0.
}
\]

That is the \textbf{Euler--Lagrange equation}.

\begin{quote}
Optimizing an integral over curves has produced a differential
equation for the curve.
\end{quote}

Two details to emphasize:
\begin{enumerate}
\item The outer \(\dd/\dd x\) is a \textbf{total derivative}
along the curve.
\item Euler--Lagrange gives a condition for \textbf{stationarity}.
It does not, by itself, prove a minimum.
\end{enumerate}

\subsection*{Checking the total derivative explicitly}

Writing the integrand as \(F(x,y,p)\), with \(p=y'\),
\[
\frac{\dd}{\dd x}F_p(x,y(x),y'(x))
=
F_{px}+F_{py}y'+F_{pp}y''.
\]

Thus the equation can also be written
\[
F_y-F_{px}-F_{py}y'-F_{pp}y''=0.
\]

This explains why Euler--Lagrange is typically a second-order
differential equation.

\subsection*{Functional derivative notation}

The same calculation can be written
\[
\delta\J[y;\eta]
=
\int_a^b
\frac{\delta\J}{\delta y(x)}
\,\eta(x)\,\dd x,
\]
where
\[
\boxed{
\frac{\delta\J}{\delta y(x)}
=
F_y-\frac{\dd}{\dd x}F_{y'}.
}
\]

The first variation is a number for each chosen direction \(\eta\).
The functional derivative above is its coefficient as a function
of \(x\).

\section*{30--39 minutes: A complete example}

Consider
\[
\boxed{
\J[y]
=
\int_0^1
\left[
\frac12(y')^2-y
\right]\dd x,
\qquad
y(0)=y(1)=0.
}
\]

Think of this as a normalized model of a stretched string under
a uniform transverse load, with \(y\) measured in the load's direction.

\begin{itemize}
\item The term \(\tfrac12(y')^2\) penalizes steep slopes.
\item The term \(-y\) favors displacement in the direction of the load.
\end{itemize}

Ask:
\begin{quote}
One term wants the string flat; the other wants it displaced.
What shape balances them?
\end{quote}

Here,
\[
F(x,y,y')=\frac12(y')^2-y.
\]

Compute each piece:
\[
F_y=-1,
\qquad
F_{y'}=y',
\qquad
\frac{\dd}{\dd x}F_{y'}=y''.
\]

Euler--Lagrange gives
\[
-1-y''=0,
\]
or
\[
y''=-1.
\]

Integrate once:
\[
y'=-x+C_1.
\]

Integrate again:
\[
y=-\frac{x^2}{2}+C_1x+C_2.
\]

Apply the endpoint conditions:
\[
y(0)=0
\quad\Longrightarrow\quad
C_2=0,
\]
and
\[
y(1)=0
\quad\Longrightarrow\quad
-\frac12+C_1=0
\quad\Longrightarrow\quad
C_1=\frac12.
\]

Thus the stationary curve is
\[
\boxed{
y_*(x)=\frac12x(1-x).
}
\]

\subsection*{Have we proved a minimum?}

Ask:
\begin{quote}
Have we proved that this is a minimum?
\end{quote}

Not yet. But for this example we can check it exactly.

Write any competing curve with the same endpoints as
\[
y=y_*+h,
\qquad
h(0)=h(1)=0.
\]

Expand the functional difference:
\begin{align*}
\J[y_*+h]-\J[y_*]
&=
\int_0^1
\left[
\frac12(y_*'+h')^2-(y_*+h)
-\frac12(y_*')^2+y_*
\right]\dd x\\
&=
\int_0^1
\left[
y_*'h'+\frac12(h')^2-h
\right]\dd x.
\end{align*}

Integrate the cross term by parts:
\begin{align*}
\int_0^1 y_*'h'\,\dd x
&=
[y_*'h]_0^1-\int_0^1 y_*''h\,\dd x\\
&=
0-\int_0^1(-1)h\,\dd x\\
&=
\int_0^1h\,\dd x.
\end{align*}

The terms linear in \(h\) cancel, leaving
\[
\boxed{
\J[y_*+h]-\J[y_*]
=
\frac12\int_0^1(h')^2\,\dd x
\geq0.
}
\]

For smooth \(h\), equality requires \(h'=0\), so \(h\) is constant.
Its endpoint values force \(h=0\).

Therefore, \(y_*\) is the \textbf{unique global minimum}
within this class of curves.

This calculation also makes the first variation visible:
at the stationary curve, the change has no linear term in the
perturbation.

\section*{39--43 minutes: Newton's law falls out}

Replace the independent variable \(x\) by time \(t\), and the unknown
function \(y\) by a particle trajectory \(q(t)\).

Consider the action
\[
S[q]
=
\int_{t_0}^{t_1}
L(q,\dot q)\,\dd t,
\]
where
\[
L(q,\dot q)=\frac12m\dot q^{\,2}-V(q).
\]

Fix the initial and final positions, and the endpoint times.

Euler--Lagrange becomes
\[
\frac{\partial L}{\partial q}
-
\frac{\dd}{\dd t}
\left(
\frac{\partial L}{\partial\dot q}
\right)
=0.
\]

Calculate:
\[
\frac{\partial L}{\partial q}=-V'(q),
\]
\[
\frac{\partial L}{\partial\dot q}=m\dot q,
\]
and
\[
\frac{\dd}{\dd t}
\left(
\frac{\partial L}{\partial\dot q}
\right)
=m\ddot q.
\]

Therefore,
\[
-V'(q)-m\ddot q=0,
\]
or
\[
\boxed{
m\ddot q=-V'(q).
}
\]

That is Newton's second law for a conservative force.

For a spring,
\[
V(q)=\frac12kq^2,
\]
so
\[
V'(q)=kq
\quad\Longrightarrow\quad
\boxed{m\ddot q+kq=0.}
\]

\begin{quote}
The integral statement describes the whole trajectory.
Requiring stationarity under arbitrary local wiggles produces
the equation of motion at every instant.
\end{quote}

Use the term \textbf{stationary action}:
a physical trajectory need not minimize the action.

\section*{43--45 minutes: Three questions to finish}

\begin{enumerate}
\item \textbf{Why must \(\eta\) vanish at the endpoints?}

Because the endpoint values of the curve are prescribed.

\item \textbf{Where does the minus sign in Euler--Lagrange come from?}

Integration by parts.

\item \textbf{Why do we need every admissible \(\eta\)?}

So stationarity of the integral forces its coefficient
to vanish pointwise.
\end{enumerate}

Leave this chain on the board:
\[
\boxed{
\begin{gathered}
y+\varepsilon\eta
\;\longrightarrow\;
\left.\frac{\dd}{\dd\varepsilon}\right|_{\varepsilon=0}
\;\longrightarrow\;
\text{integration by parts}\\[4pt]
\longrightarrow\;
F_y-\frac{\dd}{\dd x}F_{y'}=0.
\end{gathered}
}
\]

\clearpage

\section*{Instructor backup: Exact calculation for a family of wiggles}

Use the stationary curve and the variation
\[
y_*(x)=\frac12x(1-x),
\qquad
\eta(x)=x(1-x).
\]

Then
\[
y_\varepsilon(x)
=
\left(\frac12+\varepsilon\right)x(1-x).
\]

Let
\[
A=\frac12+\varepsilon.
\]
Thus,
\[
y_\varepsilon'=A(1-2x).
\]

The required integrals are
\begin{align*}
\int_0^1(1-2x)^2\,\dd x
&=
\int_0^1(1-4x+4x^2)\,\dd x\\
&=
1-2+\frac43\\
&=
\frac13,
\end{align*}
and
\[
\int_0^1x(1-x)\,\dd x
=
\frac12-\frac13
=
\frac16.
\]

Therefore,
\begin{align*}
\J[y_\varepsilon]
&=
\frac{A^2}{2}\cdot\frac13-A\cdot\frac16\\
&=
\frac{A^2-A}{6}\\
&=
\boxed{-\frac1{24}+\frac{\varepsilon^2}{6}}.
\end{align*}

Hence,
\[
\boxed{
\frac{\dd}{\dd\varepsilon}\J[y_\varepsilon]
=
\frac{\varepsilon}{3},
\qquad
\frac{\dd^2}{\dd\varepsilon^2}\J[y_\varepsilon]
=
\frac13.
}
\]

At \(\varepsilon=0\), the first derivative vanishes and
\[
\J[y_*]=-\frac1{24}.
\]

At \(\varepsilon=-\tfrac12\), the curve is flat:
\[
y_\varepsilon(x)=0,
\qquad
\J[y_\varepsilon]=0.
\]

This explores one family of curves.
The earlier calculation with arbitrary \(h\) proves the minimum
among all admissible curves.

\section*{Instructor backup: How the brachistochrone becomes a cycloid}

This is additional material, outside the 45-minute plan.

The constant factor \(1/\sqrt{2g}\) does not affect which curves
are stationary, so use
\[
F(y,y')=\frac{\sqrt{1+y'^2}}{\sqrt y}.
\]

Because \(F\) has no explicit dependence on \(x\),
Euler--Lagrange gives a convenient first integral.

\subsection*{Deriving the Beltrami identity}

Differentiate:
\begin{align*}
\frac{\dd}{\dd x}\left(F-y'F_{y'}\right)
&=
F_x+F_y y'+F_{y'}y''
-y''F_{y'}-y'\frac{\dd}{\dd x}F_{y'}\\
&=
F_x+y'\left(F_y-\frac{\dd}{\dd x}F_{y'}\right).
\end{align*}

Both terms vanish along an Euler--Lagrange solution here, so
\[
\boxed{F-y'F_{y'}=C.}
\]

This is the \textbf{Beltrami identity}.

For the present integrand,
\[
F_{y'}
=
\frac{y'}{\sqrt y\sqrt{1+y'^2}}.
\]

Therefore,
\begin{align*}
F-y'F_{y'}
&=
\frac{\sqrt{1+y'^2}}{\sqrt y}
-
\frac{y'^2}{\sqrt y\sqrt{1+y'^2}}\\
&=
\frac{1+y'^2-y'^2}
{\sqrt y\sqrt{1+y'^2}}\\
&=
\frac1{\sqrt{y(1+y'^2)}}.
\end{align*}

Consequently,
\[
y(1+y'^2)=\frac1{C^2}.
\]

Write the positive constant as \(2R\):
\[
\boxed{y(1+y'^2)=2R.}
\]

\subsection*{Solving the first integral}

Set
\[
y=R(1-\cos\theta)
=
2R\sin^2\frac\theta2.
\]

Then
\[
1+y'^2
=
\frac{2R}{y}
=
\csc^2\frac\theta2,
\]
so
\[
y'^2=\cot^2\frac\theta2.
\]

On the initial descending branch, take
\[
y'=\cot\frac\theta2.
\]

This continues smoothly through the lowest point onto the
rising branch when needed.

Since
\[
\frac{\dd y}{\dd\theta}=R\sin\theta,
\]
we obtain
\begin{align*}
\frac{\dd x}{\dd\theta}
&=
\frac{\dd y/\dd\theta}{\dd y/\dd x}\\
&=
\frac{R\sin\theta}{\cot(\theta/2)}\\
&=
R(1-\cos\theta).
\end{align*}

At the lowest point, interpret the intermediate quotient by
continuity.

Integrating and setting \(x(0)=0\),
\[
\boxed{
x(\theta)=R(\theta-\sin\theta),
\qquad
y(\theta)=R(1-\cos\theta).
}
\]

These are the parametric equations of a cycloid, with \(y\)
measured downward.

The endpoint determines \(R\) and \(\theta_B\):
\[
b=R(\theta_B-\sin\theta_B),
\qquad
h=R(1-\cos\theta_B),
\]
using
\[
0<\theta_B<2\pi.
\]

Equivalently, first solve
\[
\frac bh
=
\frac{\theta_B-\sin\theta_B}{1-\cos\theta_B},
\]
then obtain
\[
R=\frac{h}{1-\cos\theta_B}.
\]

\subsection*{Checking the travel time}

Along this arc,
\begin{align*}
\frac{\dd s}{\dd\theta}
&=
\sqrt{
\left(\frac{\dd x}{\dd\theta}\right)^2+
\left(\frac{\dd y}{\dd\theta}\right)^2
}\\
&=
R\sqrt{(1-\cos\theta)^2+\sin^2\theta}\\
&=
R\sqrt{2-2\cos\theta}\\
&=
2R\sin\frac\theta2.
\end{align*}

The speed is
\[
v=\sqrt{2gy}
=
2\sqrt{gR}\sin\frac\theta2.
\]

Thus,
\[
\dd t
=
\frac{\dd s}{v}
=
\sqrt{\frac Rg}\,\dd\theta,
\]
and
\[
\boxed{
T=\theta_B\sqrt{\frac Rg}.
}
\]

The ratio has a finite limit at release, despite both
\(\dd s/\dd\theta\) and \(v\) vanishing there.

\textit{Scope of the calculation.}
Apply the smooth Euler--Lagrange reasoning away from the singular
starting point, for example using variations supported in the
interior. This calculation identifies the cycloidal stationary
candidate. Stationarity alone does not establish global optimality.

\section*{Instructor backup: Direct Euler--Lagrange check for the brachistochrone}

For an explicit algebra check, write \(p=y'\). Then
\[
F(y,p)=\frac{\sqrt{1+p^2}}{\sqrt y},
\]
with
\[
F_y=-\frac{\sqrt{1+p^2}}{2y^{3/2}},
\qquad
F_p=\frac{p}{\sqrt y\sqrt{1+p^2}}.
\]

Differentiating \(F_p\) along the curve gives
\[
\frac{\dd}{\dd x}F_p
=
\frac{y''}{\sqrt y(1+y'^2)^{3/2}}
-
\frac{y'^2}{2y^{3/2}\sqrt{1+y'^2}}.
\]

Therefore,
\begin{align*}
F_y-\frac{\dd}{\dd x}F_p
&=
-\frac{\sqrt{1+y'^2}}{2y^{3/2}}
-\frac{y''}{\sqrt y(1+y'^2)^{3/2}}
+\frac{y'^2}{2y^{3/2}\sqrt{1+y'^2}}\\
&=
-\frac{1+y'^2+2yy''}
{2y^{3/2}(1+y'^2)^{3/2}}.
\end{align*}

Euler--Lagrange is thus equivalent, where \(y>0\), to
\[
\boxed{2yy''+1+y'^2=0.}
\]

The first integral follows directly:
\begin{align*}
\frac{\dd}{\dd x}\left[y(1+y'^2)\right]
&=
y'(1+y'^2)+2yy'y''\\
&=
y'\left(1+y'^2+2yy''\right)\\
&=
0.
\end{align*}

Hence,
\[
y(1+y'^2)=\text{constant},
\]
in agreement with the Beltrami calculation.

\section*{Small details worth checking before teaching}

\begin{itemize}
\item Fixed endpoint values imply
\[
\eta(a)=\eta(b)=0.
\]
They do not imply
\[
\eta'(a)=\eta'(b)=0.
\]

\item \(F_y\) and \(F_{y'}\) are partial derivatives;
\(\dd F_{y'}/\dd x\) is a total derivative.

\item You cannot set \(F_y=0\) and \(F_{y'}=0\) separately:
\(\eta\) and \(\eta'\) are linked.

\item In the general derivation, assuming \(F\) and the candidate
curve are sufficiently smooth justifies differentiation under
the integral and integration by parts.

\item With fixed \(x\), variation and differentiation commute:
\[
\delta y
=
\left.
\frac{\dd y_\varepsilon}{\dd\varepsilon}
\right|_{\varepsilon=0}
=
\eta,
\]
and
\[
\delta y'
=
\eta'
=
\frac{\dd}{\dd x}(\delta y).
\]

\item A stationary curve need not be a minimum.
For the string example, the exact nonnegative functional
difference supplies the additional proof.
\end{itemize}

\section*{Historical reference}

J.~J.~O'Connor and E.~F.~Robertson,
\emph{The brachistochrone problem},
MacTutor History of Mathematics, University of St Andrews.

\url{https://mathshistory.st-andrews.ac.uk/HistTopics/Brachistochrone/}

\end{document}
